\chapter{Epistemic Democracy}
\label{ch:epistemic-democracy}

\epigraph{The cure for the ailments of democracy is more democracy.}{John Dewey, \textit{The Public and Its Problems}}

\noindent
If concentrated doxonomic power distorts public belief, the remedy is the democratization of belief production.
This chapter explores what epistemic democracy means, whether it is achievable, and what institutional designs might support it.

\section{The Problem}
\label{sec:ch32-problem}

\begin{definition}[Epistemic Democracy]
\label{def:epistemic-democracy}
\index{epistemic democracy}
An \emph{epistemic democracy} is a social arrangement in which the production, circulation, and evaluation of public belief is not dominated by a narrow set of funders or institutions.
Formally, $\ehhi < \bar{h}$ for a low threshold $\bar{h}$, and the narrative infrastructure $\narr$ is responsive to diverse inputs.
\end{definition}

\section{Mechanism Design for Belief Production}
\label{sec:ch32-mechanism}

\begin{model}[Democratic Belief Production]
\label{mod:democratic-belief}
\index{mechanism design}
Treating belief production as a mechanism design problem:
\begin{itemize}[nosep]
  \item agents have private preferences over beliefs and narratives,
  \item a social planner seeks to design institutions that aggregate preferences without allowing wealthy agents to dominate,
  \item the mechanism must be incentive-compatible (agents cannot profit by misrepresenting preferences) and individually rational.
\end{itemize}
The impossibility results of \textcite{arrow1963} apply: no perfect mechanism exists.
But approximate mechanisms---public media, deliberative assemblies, open curricula---can reduce epistemic concentration.
\end{model}

\section{Institutional Designs}
\label{sec:ch32-designs}

\begin{itemize}[nosep]
  \item \textbf{Public-interest media}: funded by taxation, independent of advertising, with editorial autonomy \autocite{mcchesney1999}.
  \item \textbf{Open curricula}: educational content developed through transparent, participatory processes.
  \item \textbf{Institutional pluralism}: ensuring that no single funding source dominates any narrative domain.
  \item \textbf{Transparency requirements}: mandatory disclosure of funding for all narrative-producing institutions.
\end{itemize}

\section{Notes and References}
\label{sec:ch32-notes}

On epistemic democracy, see \textcite{dewey1927} and \textcite{habermas1962}.
Public media is analyzed in \textcite{mcchesney1999}.
Mechanism design draws on \textcite{arrow1963}.
On participatory governance, see \textcite{ostrom1990} and \textcite{wright2010}.
For democratic theory and agonistic pluralism, see \textcite{mouffe2005}.

\section{Exercises}

\begin{exercise}
Design a transparency regime for think-tank funding.
What information should be disclosed?
How would this affect the $\ehhi$?
\end{exercise}

\begin{exercise}
Compare the epistemic concentration ($\ehhi$) of a country with strong public media (e.g., BBC model) vs.\ one with purely commercial media.
What are the doxonomic consequences?
\end{exercise}

\begin{exercise}
Formalize the trade-off between epistemic diversity (high $H$) and social coordination (shared beliefs are needed for collective action).
Is there an optimal level of narrative entropy?
\end{exercise}
